% !PRINTED 347
\weberpage{347}
%<<CONT
method of approximation must lead to it. One now obtains the equations whose roots are the second, fourth, eighth \ldots\ power of equation (4):
\[
\begin{aligned}
x^3 - 4x^2 + 4x - 4 &= 0 \\
x^3 - 8x^2 - 16x - 16 &= 0 \\
x^3 - 96x^2 - 16^2 &= 0
\end{aligned}
\]
and
\[
x = \sqrt[8]{96} = 1{,}7692\ldots
\]
is already a value accurate in the first four decimals. The next step would yield no different result (because the first power of the unknown is missing in the last equation); but the step after that yields the value, accurate in the sixth decimal,
\[
x = \sqrt[32]{85032960} = 1{,}769293\ldots
\]

We shall consider yet a second example of an equation of the fifth degree, one which offers the opportunity to apply several of the theorems of the preceding section. The example is taken, like the earlier ones, from the theory of the complex multiplication of the elliptic functions. It is the equation of the fifth degree
\weq{5}{x^5 - x^3 - 2x^2 - 2x - 1 = 0.}
From this we derive the equation whose roots are the squares of the roots of (5), by substituting, in
\[
(x^5 - x^3 - 2x)^2 - (2x^2 - 1)^2 = 0
\]
$x$ for $x^2$. This gives
\weq{6}{x^5 - 2x^4 - 3x^3 - 1 = 0,}
and if we now apply the same procedure twice more, we obtain the equations whose roots are the fourth and eighth powers of the roots of (5):
\weq{7}{x^5 - 10x^4 + 9x^3 - 4x^2 - 1 = 0,}
\weq{8}{x^5 - 82x^4 + x^3 - 36x^2 - 8x - 1 = 0.}
The last of these equations lends itself to discussion. By the \gesperrt{Cartesian} theorem (\S.\,101) it has at least one positive root.

If, for equation (8), we form the functions occurring in the \gesperrt{Newton} criterion (\S.\,101)
%CONT>>
